\documentclass[12pt]{article}
\usepackage[T1]{fontenc}
\usepackage[utf8]{inputenc}
\usepackage{lmodern}
\usepackage{textcomp}
\usepackage{enumitem}
\usepackage{geometry}
\geometry{margin=1in}
\begin{document}
\begin{center}
Answer Key - Computer Science - Graduate Level Assessment 4 \
\small (Answers are ordered exactly as in the question paper.)
\end{center}

\section*{Multiple Choice}
\begin{enumerate}[label=\arabic*.]
\item \textbf{Answer:} A
\\ \textit{Options:} A) 160 bits; B) 512 bits; C) 628 bits; D) 820 bits
\\ \textit{Note:} SHA-1 produces a fixed-size message digest of 160 bits, which is a fundamental characteristic of the algorithm.
\item \textbf{Answer:} D
\\ \textit{Options:} A) Opera browser; B) Firefox; C) Chrome; D) Tor browser
\\ \textit{Note:} The Tor browser is specifically designed to enable users to access the Tor network securely and anonymously, making it essential for navigating and utilizing various services within that network.
\item \textbf{Answer:} D
\\ \textit{Options:} A) Restrict what operations/data the user can access; B) Determine if the user is an attacker; C) Flag the user if he/she misbehaves; D) Determine who the user is
\\ \textit{Note:} Authentication aims to determine who the user is by verifying their identity through various methods, ensuring that access to systems and data is granted only to authorized individuals.
\item \textbf{Answer:} B
\\ \textit{Options:} A) A procedure for testing libraries or other program components for vulnerabilities; B) Whole-system testing for security flaws and bugs; C) A security-minded form of unit testing that applies early in the development process; D) All of the above
\\ \textit{Note:} The correct answer is accurate because penetration testing involves a comprehensive evaluation of an entire system's security to identify vulnerabilities and weaknesses that could be exploited by malicious actors.
\item \textbf{Answer:} B
\\ \textit{Options:} A) .Net; B) Metasploit; C) Zeus; D) Ettercap
\\ \textit{Note:} Metasploit is a widely used penetration testing framework that simplifies the exploitation of vulnerabilities through its user-friendly interface, enabling users to perform attacks with minimal technical expertise.
\end{enumerate}

\section*{True / False}
\begin{enumerate}[label=\arabic*.]
\setcounter{enumi}{5}
\item \textbf{Answer:} True
\\ \textit{Options:} A) True; B) False
\\ \textit{Note:} A Base Transceiver Station (BTS) facilitates wireless communication for mobile operators by connecting mobile devices to the network, analogous to how an Access Point connects Wi-Fi devices to a local area network.
\item \textbf{Answer:} False
\\ \textit{Options:} A) True; B) False
\\ \textit{Note:} A buffer overflow occurs when a program writes more data to a buffer than it can hold, leading to unintended memory corruption, rather than rejecting requests when a buffer is full.
\item \textbf{Answer:} True
\\ \textit{Options:} A) True; B) False
\\ \textit{Note:} The one-time pad is considered perfectly secure because it uses a key that is as long as the message, is truly random, and is used only once, ensuring that the ciphertext does not reveal any information about the plaintext without knowledge of the key.
\item \textbf{Answer:} False
\\ \textit{Options:} A) True; B) False
\\ \textit{Note:} A fixed-length MAC tag of 5 bits is too short to provide sufficient security, as it only allows for 32 possible unique tags, making it vulnerable to brute-force and collision attacks.
\item \textbf{Answer:} True
\\ \textit{Options:} A) True; B) False
\\ \textit{Note:} In the Data Encryption Standard (DES), the key schedule algorithm processes a 56-bit key to produce a series of 48-bit subkeys used in each of the 16 rounds of encryption.
\end{enumerate}

\section*{Long Form Response}
\begin{enumerate}[label=\arabic*.]
\setcounter{enumi}{10}
\item \textbf{Answer:} def get\_item(tup 1,index): | return item
\\ \textit{Note:} correct as it demonstrates the ability to define a function that retrieves an element from a tuple at a specified index, which is fundamental in understanding tuple indexing in Python.
\item \textbf{Answer:} def max\_profit(price, k): | return final\_profit[k][n-1]
\\ \textit{Note:} correct because it succinctly encapsulates the final outcome of the dynamic programming approach used to calculate the maximum profit from at most k stock transactions, where 'final\_profit[k][n-1]' represents the maximum profit attainable with k transactions at the last day of trading.
\end{enumerate}

\vfill
\noindent\textit{End of Answer Key}
\end{document}